% TOP_PROBLEM: {"id": 3378, "title": "Gao's theorem for nonabelian groups", "category_id": 1, "category": {"id": 1, "name": "number_theory", "display_name": "Number Theory", "description": "Properties of integers, prime numbers, Diophantine equations.", "slug": "number-theory", "order_index": 1, "created_at": "2026-07-31T15:26:25.670Z"}, "status": "open", "problem_number": "OPG-337", "set_id": 12}
% FIRST_POSTED: 2026-10-01
% ENTRY_KIND: statement_only
% UnsolvedMath ID 3378; source catalogue code OPG-337
% !TeX program = lualatex
% Definition and sources only; this file is not an LLM solution attempt.
\documentclass[11pt,a4paper]{article}
\usepackage[margin=25mm]{geometry}
\usepackage{fontspec}
\setmainfont{lmroman10-regular.otf}[BoldFont=lmroman10-bold.otf,
 ItalicFont=lmroman10-italic.otf,BoldItalicFont=lmroman10-bolditalic.otf]
\usepackage{amsmath,amssymb,mathtools}
\usepackage{unicode-math}
\setmathfont{latinmodern-math.otf}
\AtBeginDocument{\let\setminus\smallsetminus}
\usepackage{xurl,hyperref}
\hypersetup{colorlinks=true,urlcolor=blue}
\setcounter{secnumdepth}{0}
\setlength{\parindent}{0pt}
\setlength{\parskip}{5pt}
\setlength{\emergencystretch}{3em}
\begin{document}
\section*{Gao's theorem for nonabelian groups}\label{3378}
{\small UnsolvedMath ID 3378 · OPG-337 · Number Theory · OpenGarden}
\subsection{Definitions and mathematical statement}
Let $G$ be a finite group, written multiplicatively with identity $1$.
A sequence is a finite indexed list of elements, with repetitions
allowed. A nonempty selection of its positions is product-one if
the selected elements can be ordered so that their product is $1$.
This ordering may differ from their order in the original list.

Define $s(G)$ to be the least integer $m$ such that every length-$m$
sequence has a product-one selection, and define $s'(G)$ in the same
way but requiring exactly $|G|$ selected positions. Both thresholds
exist: a sufficiently long sequence contains $|G|$ copies of one
element, whose product is $1$. The conjecture is
\[
                         s'(G)=s(G)+|G|-1
                   \qquad\text{for every finite group }G.
\]
These are thresholds for sequences, not sets of distinct elements.
The integer $s(G)-1$ is the largest length of a product-one-free
sequence. Appending $|G|-1$ identities to such a sequence shows
$s'(G)\ge s(G)+|G|-1$: a product-one selection of size $|G|$ would
have to use a nonempty selection of the original terms. Thus the
matching upper bound is the substantive universal claim. Noncommuting
elements are allowed throughout, and order may be chosen separately
for each candidate selection.
\subsection{Short English statement}
Does the usual relation between the two product-one subsequence thresholds hold for every finite group, including noncommutative groups?
\subsection{Sources}
\begin{itemize}
\item[S1] ULAM.AI and the UnsolvedMath Contributors, supplied problems.json, record 3378 (OPG-337), OpenGarden collection. Catalogue identity and source transcription; CC BY 4.0.
\url{https://huggingface.co/datasets/ulamai/UnsolvedMath}.
\item[S2] Open Problem Garden, Gao's theorem for nonabelian groups. Original problem and discussion; the mathematical formulation below is expanded from the supplied source record.
\url{http://www.openproblemgarden.org/op/gaos_theorem_for_nonabelian_groups}.
\end{itemize}
\end{document}
